% =====================================================================================
% APPENDICES A - K
% The bodies of A, B, G and H are written by tools/build_report_data.py directly from the
% repository, so their counts and file lists cannot drift from the artefacts they describe.
% =====================================================================================

\chapter{Repository Structure}\label{app:tree}

The repository at commit \repohead{} (\repohashfull). Directory listing to two levels; generated
output, build artefacts and the model-checker's trace files are omitted.

\begin{longtable}{p{0.40\textwidth} r p{0.42\textwidth}}
\toprule
\textbf{Directory} & \textbf{Files} & \textbf{Contents} \\
\midrule
\endhead
\texttt{(root)} & 8 & Licence, citation metadata, top-level README, reproduction guide, changelog \\
\texttt{.github/workflows} & 1 & Continuous integration \\
\texttt{FINAL\_SUBMISSION\_PACKAGE} & 2 &  \\
\texttt{FINAL\_SUBMISSION\_PACKAGE/figures} & 2 &  \\
\texttt{FINAL\_SUBMISSION\_PACKAGE/manuscript} & 2 &  \\
\texttt{FINAL\_SUBMISSION\_PACKAGE/manuscript/submission\_files} & 2 &  \\
\texttt{FINAL\_SUBMISSION\_PACKAGE/manuscript/submission\_files/similarity} & 2 &  \\
\texttt{FINAL\_SUBMISSION\_PACKAGE/reports} & 13 &  \\
\texttt{FINAL\_SUBMISSION\_PACKAGE/supplementary\_material} & 1 &  \\
\texttt{docs} & 36 & Analysis documents and the figure sources \\
\texttt{docs/figures} & 1 & The diagram generator \\
\texttt{docs/images} & 39 & Generated SVG, PDF and PNG figures \\
\texttt{docs/media} & 2 &  \\
\texttt{experiments} & 17 & Harnesses, analysis scripts, figure generators, verification checks \\
\texttt{experiments/broker} & 3 & Real Mosquitto testbed: setup, ACL variants, scenario runner \\
\texttt{experiments/corpus} & 8 & Corpus builder, static rule, controls, commit pins, findings \\
\texttt{experiments/hardware} & 4 & Probe firmware, collection and pre-registered analysis (not yet run) \\
\texttt{experiments/hardware/latency\_probe} & 1 &  \\
\texttt{experiments/results} & 15 & Committed raw and aggregate experimental records \\
\texttt{firmware/espnow\_receiver} & 3 & Companion receiver sketch (not part of the evaluation) \\
\texttt{firmware/smart\_home} & 6 & The artefact under study: v2.2.0 sketch and its headers \\
\texttt{formal} & 7 & TLA+ specifications, three TLC configurations, runner, tla2tools.jar \\
\texttt{paper} & 7 & Manuscript source and build products \\
\texttt{paper/submission} & 2 & Highlights, declaration document, similarity-screening file \\
\texttt{paper/submission/similarity} & 2 &  \\
\texttt{tests/host} & 2 & Host test suite and the stub headers it compiles against \\
\texttt{tests/host/stubs} & 9 &  \\
\texttt{tests/host/stubs/freertos} & 4 &  \\
\texttt{tools} & 5 & Command signer, evaluator, figure-package builders \\
\texttt{wokwi} & 2 & Simulator circuit definition \\
\bottomrule
\end{longtable}


\chapter{Source Code Inventory}\label{app:inventory}

Every source file in the repository with its language, total lines and non-comment,
non-blank lines. \textbf{\srcfiles{} files, \srccode{} lines of code.} The firmware itself is a
minority of the total; the verification and experimental infrastructure is the larger part, which
is a deliberate consequence of treating that infrastructure as part of the artefact
(Chapter~\ref{ch:tools}).

The complete source is distributed as
\code{FINAL\_SUBMISSION\_PACKAGE/source\_code/COMPLETE\_SOURCE\_CODE\_AND\_EXPERIMENTS.zip}, which
preserves the directory structure and contains no credentials, keys or environment-specific
files. The repository's only key material is the published demonstration string in
\code{secrets.example.h}, which the build explicitly refuses to accept as a real key.

\begin{longtable}{p{0.50\textwidth} l r r}
\toprule
\textbf{File} & \textbf{Language} & \textbf{Lines} & \textbf{Code} \\
\midrule
\endhead
\texttt{\footnotesize docs/figures/make\_diagrams.py} & Python & 1213 & 971 \\
\texttt{\footnotesize experiments/ablation.cpp} & C++ & 632 & 384 \\
\texttt{\footnotesize experiments/analyse.py} & Python & 172 & 141 \\
\texttt{\footnotesize experiments/analyse\_ablation.py} & Python & 377 & 333 \\
\texttt{\footnotesize experiments/broker/run.sh} & Shell & 32 & 14 \\
\texttt{\footnotesize experiments/broker/setup.sh} & Shell & 93 & 51 \\
\texttt{\footnotesize experiments/broker/validate\_acl.py} & Python & 287 & 239 \\
\texttt{\footnotesize experiments/campaign.cpp} & C++ & 289 & 209 \\
\texttt{\footnotesize experiments/check\_figure\_fit.py} & Python & 135 & 103 \\
\texttt{\footnotesize experiments/check\_figure\_type.py} & Python & 71 & 56 \\
\texttt{\footnotesize experiments/corpus/analyse\_guards.py} & Python & 218 & 176 \\
\texttt{\footnotesize experiments/corpus/build\_corpus.py} & Python & 136 & 119 \\
\texttt{\footnotesize experiments/corpus/fetch\_corpus.sh} & Shell & 42 & 28 \\
\texttt{\footnotesize experiments/corpus/test\_rule.py} & Python & 99 & 79 \\
\texttt{\footnotesize experiments/figstyle.py} & Python & 68 & 51 \\
\texttt{\footnotesize experiments/fuzz\_cmd.cpp} & C++ & 173 & 129 \\
\texttt{\footnotesize experiments/hardware/analyse\_hw.py} & Python & 110 & 92 \\
\texttt{\footnotesize experiments/hardware/collect.py} & Python & 65 & 55 \\
\texttt{\footnotesize experiments/hardware/latency\_probe/latency\_probe.ino} & Arduino/C++ & 53 & 31 \\
\texttt{\footnotesize experiments/make\_ablation\_figures.py} & Python & 182 & 138 \\
\texttt{\footnotesize experiments/make\_figures.py} & Python & 91 & 74 \\
\texttt{\footnotesize experiments/run.sh} & Shell & 32 & 21 \\
\texttt{\footnotesize experiments/run\_ablation.sh} & Shell & 50 & 35 \\
\texttt{\footnotesize experiments/sensitivity.cpp} & C++ & 145 & 86 \\
\texttt{\footnotesize experiments/test\_conformance.py} & Python & 165 & 132 \\
\texttt{\footnotesize experiments/test\_harness.py} & Python & 174 & 141 \\
\texttt{\footnotesize experiments/validate\_figures.py} & Python & 153 & 126 \\
\texttt{\footnotesize firmware/espnow\_receiver/alertmsg.h} & C/C++ header & 78 & 57 \\
\texttt{\footnotesize firmware/espnow\_receiver/authmac.h} & C/C++ header & 132 & 109 \\
\texttt{\footnotesize firmware/espnow\_receiver/espnow\_receiver.ino} & Arduino/C++ & 80 & 55 \\
\texttt{\footnotesize firmware/smart\_home/alertmsg.h} & C/C++ header & 78 & 57 \\
\texttt{\footnotesize firmware/smart\_home/authmac.h} & C/C++ header & 132 & 109 \\
\texttt{\footnotesize firmware/smart\_home/secrets.example.h} & C/C++ header & 43 & 15 \\
\texttt{\footnotesize firmware/smart\_home/sensor\_fusion.h} & C/C++ header & 58 & 37 \\
\texttt{\footnotesize firmware/smart\_home/smart\_home.ino} & Arduino/C++ & 1166 & 836 \\
\texttt{\footnotesize firmware/smart\_home/wifi\_guard.h} & C/C++ header & 61 & 41 \\
\texttt{\footnotesize formal/AlarmFSM.cfg} & TLC config & 14 & 10 \\
\texttt{\footnotesize formal/AlarmFSM.tla} & TLA+ & 146 & 124 \\
\texttt{\footnotesize formal/AlarmFSM\_unsafe.cfg} & TLC config & 12 & 7 \\
\texttt{\footnotesize formal/Isolation.cfg} & TLC config & 14 & 10 \\
\texttt{\footnotesize formal/Isolation.tla} & TLA+ & 143 & 124 \\
\texttt{\footnotesize formal/run\_tlc.sh} & Shell & 48 & 34 \\
\texttt{\footnotesize tests/host/run\_tests.sh} & Shell & 78 & 63 \\
\texttt{\footnotesize tests/host/stubs/Adafruit\_GFX.h} & C/C++ header & 4 & 4 \\
\texttt{\footnotesize tests/host/stubs/Adafruit\_SSD1306.h} & C/C++ header & 16 & 16 \\
\texttt{\footnotesize tests/host/stubs/Arduino.h} & C/C++ header & 61 & 61 \\
\texttt{\footnotesize tests/host/stubs/Preferences.h} & C/C++ header & 16 & 15 \\
\texttt{\footnotesize tests/host/stubs/PubSubClient.h} & C/C++ header & 33 & 33 \\
\texttt{\footnotesize tests/host/stubs/WiFi.h} & C/C++ header & 18 & 18 \\
\texttt{\footnotesize tests/host/stubs/WiFiClientSecure.h} & C/C++ header & 4 & 4 \\
\texttt{\footnotesize tests/host/stubs/Wire.h} & C/C++ header & 4 & 4 \\
\texttt{\footnotesize tests/host/stubs/esp\_system.h} & C/C++ header & 7 & 7 \\
\texttt{\footnotesize tests/host/stubs/freertos/FreeRTOS.h} & C/C++ header & 29 & 29 \\
\texttt{\footnotesize tests/host/stubs/freertos/queue.h} & C/C++ header & 1 & 1 \\
\texttt{\footnotesize tests/host/stubs/freertos/semphr.h} & C/C++ header & 1 & 1 \\
\texttt{\footnotesize tests/host/stubs/freertos/task.h} & C/C++ header & 1 & 1 \\
\texttt{\footnotesize tests/host/test\_firmware.cpp} & C++ & 510 & 412 \\
\texttt{\footnotesize tools/build\_figure\_collection.py} & Python & 191 & 155 \\
\texttt{\footnotesize tools/build\_figure\_preview.py} & Python & 142 & 125 \\
\texttt{\footnotesize tools/build\_report\_data.py} & Python & 377 & 315 \\
\texttt{\footnotesize tools/build\_source\_archive.py} & Python & 195 & 153 \\
\texttt{\footnotesize tools/eval\_commands.py} & Python & 125 & 91 \\
\texttt{\footnotesize tools/measure.py} & Python & 122 & 99 \\
\texttt{\footnotesize tools/sign\_command.py} & Python & 61 & 47 \\
\midrule
\textbf{64 files} & & \textbf{9458} & \textbf{7293} \\
\bottomrule
\end{longtable}


\chapter{Firmware Configuration Reference}\label{app:config}

\section{Build-time switches}

\begin{longtable}{@{}p{0.34\textwidth}p{0.10\textwidth}p{0.47\textwidth}@{}}
\toprule
\textbf{Macro} & \textbf{Default} & \textbf{Effect} \\
\midrule
\endhead
\code{AUTH\_COMMANDS}      & 0 & Requires an HMAC-SHA256 tag and a strictly increasing counter on
every remote command. \textbf{A deployment should set this to 1.} \\
\code{ALLOW\_REMOTE\_DISARM}& 0 & Permits remote disarm. Combined with
\code{AUTH\_COMMANDS~0} it is rejected at compile time unless the acknowledgement macro below is
defined. \\
\code{ACKNOWLEDGE\_UNAUTHENTICATED\_DISARM} & undefined & The deliberate escape hatch. Named for
what the builder is accepting. Exists only so the research baseline stays reproducible. \\
\code{MQTT\_USE\_TLS}      & --- & TLS transport to the broker. \\
\code{SECRET\_CMD\_KEY}    & demo & A public demonstration string. A build enabling
\code{AUTH\_COMMANDS} without overriding it is rejected at compile time, because authentication
against a published key is not authentication. \\
\bottomrule
\end{longtable}

\section{Compile-time constants}

\begin{longtable}{@{}p{0.30\textwidth}p{0.12\textwidth}p{0.49\textwidth}@{}}
\toprule
\textbf{Constant} & \textbf{Value} & \textbf{Why this value} \\
\midrule
\endhead
\code{CMD\_Q\_LEN}            & 4        & Command queue depth. Two \code{static\_assert}s tie it
to the drain bound. \\
\code{CMD\_DRAIN\_PER\_PASS}  & 4        & Equal to \code{CMD\_Q\_LEN}, which is what makes the
loop's execution time statically bounded (Section~\ref{sec:impl-drain}). \\
Event queue depth             & 16       & Drop-oldest. The queue never blocks the safety task.
\\
\code{AUTH\_TAG\_BYTES}       & 16       & 128-bit truncated HMAC-SHA256 tag. \\
\code{DEAUTH\_THRESHOLD}      & 12       & Raised from 8 after the sensitivity sweep showed the
original value firing in 79.9\% of benign runs. \\
Detection window              & 2000\,ms & Held fixed while the threshold was swept. \\
Exit / entry delay            & 10\,s    & \\
Alarm duration                & 30\,s    & \\
\bottomrule
\end{longtable}

\section{Command wire format}

\begin{lstlisting}[style=sh,caption={The authenticated command format. The counter is committed to
non-volatile storage \emph{before} the command is acted upon.}]
VERB:CTR:hex(HMAC-SHA256(key, "VERB:CTR")[0:16])
\end{lstlisting}

Committing the counter before acting is deliberate: a reset between acting and committing would
otherwise leave the counter below a value already used, reopening the replay window.

\chapter{Formal Specification Listings}\label{app:formal}

The complete text of \code{formal/AlarmFSM.tla} and \code{formal/Isolation.tla} is in the
repository and in the source archive. This appendix records the configurations and their required
outcomes; Chapter~\ref{ch:formal} discusses them.

\begin{longtable}{@{}p{0.23\textwidth}p{0.26\textwidth}p{0.20\textwidth}p{0.22\textwidth}@{}}
\toprule
\textbf{Configuration} & \textbf{Constants} & \textbf{Result} & \textbf{Required} \\
\midrule
\endhead
\code{AlarmFSM.cfg} & \code{MaxCtr~=~4}, \code{AllowUnauthDisarm~=~FALSE} & 166 generated, 25
distinct, no error & pass \\
\code{Isolation.cfg} & \code{QCap~=~16}, \code{MaxEvents~=~8} & 2,686 generated, 660 distinct, no
error & pass \\
\code{AlarmFSM\_unsafe.cfg} & \code{MaxCtr~=~4}, \code{AllowUnauthDisarm~=~TRUE} &
\code{NoUnauthRemoteDisarm} violated & \textbf{violation}; a pass fails the build \\
\bottomrule
\end{longtable}

All three are re-run by \code{formal/run\_tlc.sh} on every push. The third requires the specific
string \code{Invariant NoUnauthRemoteDisarm is violated} in TLC's output, not merely a non-zero
exit, so that a syntax error or a missing constant cannot be mistaken for the expected violation.

\chapter{Test Suite Reference}\label{app:tests}

\begin{longtable}{@{}p{0.33\textwidth}p{0.17\textwidth}p{0.42\textwidth}@{}}
\toprule
\textbf{Suite} & \textbf{Scale} & \textbf{What it would catch} \\
\midrule
\endhead
\code{tests/host/run\_tests.sh} & 129 assertions & Firmware logic defects. Includes three
\emph{negative build} tests that require a defective configuration to fail to compile; a
successful build is a test failure. \\
\code{formal/run\_tlc.sh} & 3 configurations & A vacuous invariant recurring silently. \\
\code{experiments/test\_conformance.py} & 19,529 transitions & The firmware diverging from the
specification. Includes a negative control that must be rejected. \\
\code{experiments/test\_harness.py} & metamorphic and differential properties & Harness defects:
determinism, seed sensitivity, monotone attack effort, mode separation, and agreement of the HMAC
payload format with an independent implementation. \\
\code{experiments/corpus/test\_rule.py} & 8 controls & A static rule whose polarity is wrong.
Includes this project's own pre-repair guard as a positive control and its repaired guard as a
negative control. \\
\code{experiments/validate\_figures.py} & 59 assertions & A figure drifting from the raw records,
or a baked-in cross-reference. \\
\code{experiments/check\_figure\_type.py} & all figures & A label printing below 6\,pt. \\
\code{experiments/check\_figure\_fit.py} & all labels & A clipped label, measured in the real
font. \\
\bottomrule
\end{longtable}

All suites pass at commit \repohead{}.

\chapter{Build and Deployment Instructions}\label{app:build}

\section{Prerequisites}

\code{arduino-cli} with the \code{esp32:esp32} core; the \code{Adafruit SSD1306},
\code{Adafruit GFX} and \code{PubSubClient} libraries; GCC 13 or later for the host tests; Python
3.11 or later; Java for \code{tla2tools.jar}; Mosquitto 2.0.x and OpenSSL for the broker
experiment; and \TeX{} Live for the documents.

\section{Firmware}

\begin{lstlisting}[style=sh]
cp firmware/smart_home/secrets.example.h firmware/smart_home/secrets.h
# Edit secrets.h: set Wi-Fi credentials and REPLACE the demonstration keys.
# The build refuses AUTH_COMMANDS=1 with the published demonstration key.

arduino-cli compile --fqbn esp32:esp32:esp32 firmware/smart_home \
  --build-property "compiler.cpp.extra_flags=-DAUTH_COMMANDS=1"
arduino-cli upload  --fqbn esp32:esp32:esp32 -p /dev/ttyUSB0 firmware/smart_home
\end{lstlisting}

\section{Verification and experiments}

\begin{lstlisting}[style=sh]
bash tests/host/run_tests.sh          # 129 assertions, incl. negative build tests
bash formal/run_tlc.sh                # 3 configurations; the third must report a violation
bash experiments/run.sh               # E1, E2: the adversarial campaign
bash experiments/run_ablation.sh      # E3: 12 configurations, 3,600 trials
bash experiments/broker/setup.sh      # E6: CA, certificates, password and ACL files
bash experiments/broker/run.sh        # E6: 8 scenarios against Mosquitto
python3 experiments/validate_figures.py
\end{lstlisting}

\section{Documents}

\begin{lstlisting}[style=sh]
cd paper  && pdflatex manuscript.tex   # three passes
cd report && pdflatex report.tex       # three passes, for the ToC, LoF and LoT
\end{lstlisting}

\section{Deployment guidance}

A deployment should set \code{AUTH\_COMMANDS~1} with a key that is not the published default,
use per-client broker credentials with a deny-by-default access control list that does
\emph{not} grant the device identity write access on the command topic, and enable flash
encryption and secure boot --- which this firmware does not. The two open design problems of
Section~\ref{ch:limitations} apply to any such deployment and are not fixed.

\chapter{Experiment Catalogue}\label{app:experiments}

Every experiment with its evidence class, artefact and scale. The evidence class is the
distinction maintained throughout this report: a formal result and a simulated result support
different claims, and the hardware conditions support none because they produced no data.

\begin{longtable}{l p{0.30\textwidth} p{0.15\textwidth} p{0.17\textwidth} p{0.20\textwidth}}
\toprule
\textbf{ID} & \textbf{Experiment} & \textbf{Evidence class} & \textbf{Artefact} & \textbf{Scale} \\
\midrule
\endhead
\texttt{E1} & Adversarial campaign, baseline build & Simulation & \texttt{\footnotesize summary.json} & 3{,}000 trials \\
\texttt{E2} & Adversarial campaign, hardened build & Simulation & \texttt{\footnotesize summary.json} & 3{,}000 trials \\
\texttt{E3} & Command-path ablation & Simulation & \texttt{\footnotesize ablation.jsonl} & 3{,}600 trials, 12 configurations \\
\texttt{E4} & Host-build fuzzing under ASan/UBSan & Simulation & \texttt{\footnotesize fuzz.json} & 3{,}000{,}000 iterations \\
\texttt{E5} & Detector parameter sensitivity sweep & Simulation & \texttt{\footnotesize sensitivity.jsonl} & 180 grid cells \\
\texttt{E6} & Access control against a live broker & Testbed measurement & \texttt{\footnotesize broker\_acl.jsonl} & 8 scenarios \\
\texttt{E7} & Open-source corpus study & Corpus study & \texttt{\footnotesize corpus/corpus.json} & 135 projects \\
\texttt{E8} & Trace conformance against the TLA+ relation & Formal + simulation & \texttt{\footnotesize test\_conformance.py} & 19{,}529 transitions \\
\texttt{E9} & Model checking, three configurations & Formal & \texttt{\footnotesize formal/} & 25 and 660 distinct states \\
\texttt{---} & Hardware conditions A--D & \emph{not executed} & \texttt{\footnotesize ---} & \emph{no data exists} \\
\bottomrule
\end{longtable}


\chapter{Figure Index}\label{app:figures}

Every figure with its title, file and the script that generates it. No figure was drawn by hand.

\begin{longtable}{l p{0.30\textwidth} p{0.26\textwidth} p{0.26\textwidth}}
\toprule
\textbf{No.} & \textbf{Title} & \textbf{File} & \textbf{Generator} \\
\midrule
\endhead
1 & System architecture & \texttt{\footnotesize architecture.pdf} & \texttt{\footnotesize docs/figures/make\_diagrams.py} \\
2 & Command execution flow & \texttt{\footnotesize command\_flow.pdf} & \texttt{\footnotesize docs/figures/make\_diagrams.py} \\
3 & Command-path ablation results & \texttt{\footnotesize ablation\_results.pdf} & \texttt{\footnotesize experiments/make\_ablation\_figures.py} \\
4 & Broker access control architecture & \texttt{\footnotesize acl\_architecture.pdf} & \texttt{\footnotesize docs/figures/make\_diagrams.py} \\
5 & Security compromise comparison & \texttt{\footnotesize credential\_matrix.pdf} & \texttt{\footnotesize experiments/make\_ablation\_figures.py} \\
6 & Experimental setup & \texttt{\footnotesize experimental\_setup.pdf} & \texttt{\footnotesize docs/figures/make\_diagrams.py} \\
7 & Compile-guard comparison & \texttt{\footnotesize guard\_decision.pdf} & \texttt{\footnotesize docs/figures/make\_diagrams.py} \\
8 & Transport security versus command authorisation & \texttt{\footnotesize security\_paths.pdf} & \texttt{\footnotesize docs/figures/make\_diagrams.py} \\
\midrule
\multicolumn{4}{l}{\textit{Repository documentation figures, not part of the manuscript}}\\
\midrule
--- & Command sequence (UML) & \texttt{\footnotesize command\_sequence.pdf} & \texttt{\footnotesize docs/figures/make\_diagrams.py} \\
--- & Deployment view & \texttt{\footnotesize deployment.pdf} & \texttt{\footnotesize docs/figures/make\_diagrams.py} \\
--- & Methodology overview & \texttt{\footnotesize methodology.pdf} & \texttt{\footnotesize docs/figures/make\_diagrams.py} \\
--- & Alarm state machine & \texttt{\footnotesize state\_machine.pdf} & \texttt{\footnotesize docs/figures/make\_diagrams.py} \\
--- & Event sequence & \texttt{\footnotesize event\_sequence.pdf} & \texttt{\footnotesize docs/figures/make\_diagrams.py} \\
--- & Circuit wiring & \texttt{\footnotesize circuit\_wiring.pdf} & \texttt{\footnotesize docs/figures/make\_diagrams.py} \\
--- & Hardware measurement rig (not built) & \texttt{\footnotesize hardware\_rig.pdf} & \texttt{\footnotesize docs/figures/make\_diagrams.py} \\
\bottomrule
\end{longtable}


\chapter{Hardware Validation Protocol}\label{app:hardware}

\textbf{This protocol was not executed. No hardware data exists.} It is published so that the
timing claim left open by this work is falsifiable by someone else, and so that the analysis
cannot be chosen after the data is seen.

\section{The probe that established hardware was unavailable}

Recorded so that the absence is a finding rather than an assertion. In the environment that
produced the committed results there were no \code{/dev/ttyUSB*} or \code{/dev/ttyACM*} devices,
no USB subsystem, no \code{arduino-cli} on the path, no Python \code{serial} module, and the only
network interfaces present were \code{eth0}, \code{ifb0} and \code{lo} --- no wireless interface
of any kind. An RF-shielded enclosure, a hard prerequisite for condition D, was likewise
unavailable.

\section{The rig}

A device under test and a separate probe microcontroller. The device under test raises a GPIO pin
on sensor confirmation and again on siren assertion; the probe watches both and timestamps them
with its own hardware timer, so the measurement does not depend on the device under test's own
clock or on its serial output.

\textbf{The two boards must share a common ground.} Without it they have no shared reference and
the readings are meaningless. This is the most common way this rig is got wrong.

\section{The four conditions}

\begin{longtable}{@{}p{0.09\textwidth}p{0.29\textwidth}p{0.11\textwidth}p{0.42\textwidth}@{}}
\toprule
\textbf{Cond.} & \textbf{Setup} & \textbf{Shielding} & \textbf{What it would establish} \\
\midrule
\endhead
A --- idle & Wi-Fi connected, broker reachable, no other traffic & none & The control
distribution and the probe's own noise floor. Everything else is compared against this. \\
B --- legitimate load & Frequent status publication, large retained payloads & none & Whether
ordinary operation perturbs the safety path. \textbf{A difference here would be the most damaging
result for the design}, because it needs no adversary. \\
C --- command stress & High-rate publication on the device's own command topic & none & The
hardware counterpart of the simulated flood: where the isolation claim is actually tested. \\
D --- radio stress & 802.11 deauthentication and disassociation frames & \textbf{required} &
The central open question: whether the mitigation holds under load on the radio itself rather
than on the IP path. \\
\bottomrule
\end{longtable}

\textbf{The device must not be reflashed between conditions.} The firmware must be byte-identical
across all four, or the comparison is between two systems rather than two load levels.

\section{The pre-registered analysis}

Committed before any data existed, and \textbf{not to be substituted after seeing the data}.

Two-sided Mann--Whitney $U$ test with tie correction, each condition against condition A.
$\alpha = 0.05$, power $0.80$, equal $n$ per group. The planned sample size is
\textbf{$n = 200$ per condition}, which detects a shift of 0.29 standard deviations --- 29\,\textmu s
at an idle spread of 100\,\textmu s, 52\,\textmu s at 180\,\textmu s, 87\,\textmu s at
300\,\textmu s. The effect size is expressed in standard deviations because the idle spread is
unknown until condition A is run.

\begin{longtable}{@{}rrrrrr@{}}
\toprule
$n$ per cond. & $A = P(X>Y)$ & Shift (SD) & sd $=100$\,\textmu s & sd $=180$\,\textmu s &
sd $=300$\,\textmu s \\
\midrule
\endhead
30  & 0.709 & 0.78 & 78 & 140 & 233 \\
50  & 0.662 & 0.59 & 59 & 106 & 177 \\
100 & 0.614 & 0.41 & 41 & 74  & 123 \\
\textbf{200} & \textbf{0.581} & \textbf{0.29} & \textbf{29} & \textbf{52} & \textbf{87} \\
400 & 0.557 & 0.20 & 20 & 37  & 61 \\
\bottomrule
\end{longtable}

A non-parametric test is pre-registered rather than a $t$-test because interrupt-driven latency
distributions are routinely skewed with a long upper tail, and because the 95th percentile and
maximum matter more for a safety function than the mean.

\section{Status of the tooling}

\begin{longtable}{@{}p{0.40\textwidth}p{0.50\textwidth}@{}}
\toprule
\textbf{Component} & \textbf{Status} \\
\midrule
\endhead
\code{experiments/hardware/latency\_probe/} & Probe firmware, written, not flashed \\
\code{experiments/hardware/collect.py} & Serial collection, written, appends across conditions \\
\code{experiments/hardware/analyse\_hw.py} & Summary statistics, the pre-registered
Mann--Whitney $U$, and the box plot. \textbf{Verified end to end on explicitly synthetic input},
labelled as a script self-test. The synthetic file was deliberately not committed, so that no
reader can mistake it for data. \\
\code{results\_hw\_TEMPLATE.csv} & Schema only:
\code{condition,index,latency\_us}. \textbf{Contains no rows.} \\
\code{docs/HARDWARE\_A\_TO\_C\_RUNBOOK.md} & Step-by-step procedure for A--C, executable by
someone else \\
\bottomrule
\end{longtable}

Conditions A--C require an ESP32 DevKit, a logic analyser or a second microcontroller as a timing
reference, and a controlled network. Condition D additionally requires an RF-shielded enclosure
and the appropriate permission; the repository publishes no attack tooling for it.

\chapter{Journal Submission Status}\label{app:journal}

Recorded as status rather than as achievement. The manuscript has \textbf{not} been submitted to
any journal.

\begin{longtable}{@{}p{0.31\textwidth}p{0.17\textwidth}p{0.44\textwidth}@{}}
\toprule
\textbf{Requirement} & \textbf{Status} & \textbf{Evidence or action needed} \\
\midrule
\endhead
Guide for Authors obtained & verified & Retrieved from the publisher's own site \\
Page and reference limits & verified & 31 pages, 22 references, 0 dangling \\
Abstract length & verified & 249 words, within the stated limit \\
Keywords & verified & 7, within the stated limit \\
Structured declarations & verified & Present in the publisher's named sections \\
Highlights & verified & 5 bullets, longest 80 of 85 characters \\
Figure format and font embedding & verified & All 8 fully embedded, 0 Type~3 \\
Journal quartile and metric years & \textbf{unverified} & The ranking source returned an
automated challenge. \textbf{Not circumvented}; no quartile is asserted anywhere in this work. \\
Competing-interest declaration & \textbf{author confirmation required} & Drafted; must not be
declared on the authors' behalf \\
Funding declaration & \textbf{author confirmation required} & Drafted; no funding is invented \\
Authorship contribution statement & \textbf{author confirmation required} & The sole-author
versus supervisor-co-author question is the author's to settle \\
Corresponding author's email and postal address & \textbf{required} & Not present on the title
page \\
Institutional similarity report & \textbf{outstanding} & No licensed tool was available.
\textbf{No similarity percentage is claimed, official or otherwise.} \\
Supervisor approval & \textbf{outstanding} & Not obtained \\
Dissertation-overlap ruling & \textbf{outstanding} & Not made \\
Five missing DOIs & \textbf{outstanding} & \code{mcreview2026}, \code{espwn2023},
\code{mqtttls2026}, \code{zeghida2024}, and an ISBN for \code{lamport2002} \\
\bottomrule
\end{longtable}

\chapter{Reproducibility Record}\label{app:repro}

\section{What was re-verified during the preparation of this report}

\begin{itemize}[leftmargin=1.6em,itemsep=4pt]
  \item \textbf{Every reported count against its raw artefact.} 6,000 campaign trials; 3,600
        ablation trials (9,600 adversarial trials in total); 3,000,000 fuzz iterations; 180
        sensitivity cells; 8 broker scenarios; 135 corpus projects; 19,529 trace transitions.
  \item \textbf{Both Fisher $p$-values recomputed} from the raw records:
        $3.53\times10^{-184}$ and $9.83\times10^{-119}$.
  \item \textbf{The Wilson interval recomputed:} [12.9, 21.3].
  \item \textbf{All three TLC configurations re-run:} 25 and 660 distinct states, and the
        required violation in the third.
  \item \textbf{The broker experiment re-run} against a freshly installed Mosquitto 2.0.18 in a
        clean environment. All eight verdicts identical; only log timestamps and random client
        identifiers differed.
  \item \textbf{Every figure regenerated.} All SVGs byte-identical; the PDFs differ only in
        embedded creation timestamps.
  \item \textbf{The full test suite:} 129 host assertions, 8 corpus controls, the harness
        metamorphic and differential properties, and 59 figure assertions. All pass.
\end{itemize}

\section{Determinism}

Seeds are fixed at 1000, 2000 and 7. The ablation harness forks per trial, so one trial cannot
influence the next through accumulated state. All statistics are implemented with the Python
standard library only, so no external package version can alter a reported value.

\section{What cannot be reproduced}

The hardware conditions, because they were never run. The corpus study depends on the state of
third-party repositories at fetch time and on the GitHub search API, so a re-fetch will not
return the same 135 projects; the fetched corpus and the rejection counts are committed for this
reason, and the eight rule controls are reproducible independently of the corpus.
